\chapter{Основные положения}

Определение мореходных качеств судна, а именно: плавучести, остойчивости, непотопляемости, ходкости, мореходности и управляемости — связано с вычислением гидростатических и гидродинамических сил и моментов, действующих на смоченную поверхность в различных условиях эксплуатации.
Величина, направление и точки приложения этих сил зависят от формы корпуса.
Форма корпуса может быть задана аналитически, в виде таблиц и графически.

В аналитическом виде форма корпуса задается, как правило, для научно-\-иссле\-дова\-тельских целей и эталонных моделей или судов.
К этому способу описания можно отнести также использование аналитических аппроксимаций эмпирических линий формы корпуса.

\section{Теоретический чертеж и его характеристики}

При табличном описании формы корпуса судна производится табулирование отдельных характерных точек поверхности. Примером такого представления корпуса является таблица ординат.

Графически корпус судна представляется в виде теоретического чертежа, состоящего из трех основных проекций.

В практике проектирования наибольшее распространение получило изображение формы корпуса судна в виде теоретического чертежа с описанием поверхности в табличной форме.
Отсутствие аналитических выражений для формы корпуса в целом и отдельных его кривых (шпангоутов, ватерлиний, батоксов, рыбин и т.д.) усложняет расчеты мореходных качеств судна и позволяет определить их лишь приближенно.

На теоретическом чертеже изображается поверхность судна без наружной обшивки.
Исключение составляют деревянные суда, наружная обшивка которых совпадает с поверхностью чертежа.

Теоретический чертеж судна вычерчивается в трех взаимно перпендикулярных проекциях, представляющих собой изображение корпуса на основную, диаметральную и мидель-шпангоута плоскости.
Эти проекции соответственно называются полуширотой, боком и корпусом, или основными плоскостями теоретического чертежа.
Для построения теоретического чертеж, расчетов и представления о форме корпуса судна последний рассекается дополнительными плоскостями, параллельными основным.
Кривые, образующиеся от пересечения этих плоскостей с поверхностью теоретического чертежа, носят названия: ватерлиний, батоксов и шпангоутов.
Ватерлинии образуются плоскостями, параллельными основной плоскости, батоксы — диаметральной плоскости, шпангоуты — плоскости мидель-шпангоута.
Последний находится в середине длины между перпендикулярами, под которой для двухвинтовых судов понимается расстояние между носовыми и кормовыми перпендикулярами, опущенными через точки пересечения фор- и ахтерштевней с конструктивной (КВЛ) ватерлинией.
У одновинтовых судов кормовой перпендикуляр совпадает с осью баллера руля: из-за этого мидель-шпангоут находится не в середине.
В силу симметрии корпуса судна относительно диаметральной плоскости на теоретическом чертеже показывают лишь половинки ватерлиний и шпангоутов.

При этом на правой от диаметральной плоскости части проекции изображаются ветви носовых шпангоутов, на левой — кормовых.

Количество шпангоутов, батоксов и ватерлиний определяется точностью построения чертежа и соответствующих расчетов по нему.
Так, количество ватерлиний колеблется от шести до 20, шпангоутов от 11 до 40, батоксов от двух до шести.
Первые значения относятся к упрощенным экспресс-расчетам ручным способом, последние — к вычислениям с использованием вычислительной техники.

Нумерация шпангоутов, принятая в России, начинается со шпангоута, совпадающего с носовым перпендикуляром, которому присваивается нулевой номер.
Ватерлинии нумеруются от основной плоскости, последняя получает нулевой номер.
За нулевой батокс принимается диаметральная плоскость.
Номера батоксов возрастают по мере удаления от диаметральной плоскости.

Основные размеры теоретического чертежа получили названия главных размерений.
К ним относятся:

--- длина наибольшая $L_{\text{НБ}}$ --- расстояние между перпендикулярами, опущенными из крайних носовой и кормовой точек корпуса судна;

--- длина по КВЛ $L_{\text{КВЛ}}$ --- расстояние между точками ее пересечения с фор- и ахтерштевнями;

--- длина между перпендикулярами $L_{\text{ПП}}$ --- см. выше;

--- ширина наибольшая $B_{\text{НБ}}$ --- расстояние между плоскостями, параллельными диаметральной и касательными к наружной поверхности без учета выступающих частей;

--- ширина по КВЛ $B_{\text{КВЛ}}$ --- расстояние между касательными и конструктивной ватерлинии, проведенными параллельной диаметральной плоскости;

--- высота борта $H$ --- вертикальное расстояние от основной плоскости до бортовой линии верхней палубы, измеренное в плоскости мидель-шпангоута;

--- осадка по КВЛ $T_{\text{КВЛ}}$ --- вертикальное расстояние от основной плоскости до конструктивной ватерлинии, измеренное в плоскости мидель-шпангоута.

Главные размерения наносятся на теоретическом чертеже и позволяют судить об абсолютных размерах судна.
Сами по себе они не дают представления о форме корпуса.
Некоторые сведения о ней дают соотношения главных размерений.
Так, например, отношение $L/B$ характеризует удлинение корпуса и используется в задачах ходкости, мореходности, управляемости, отношения $B/T$ и $B/H$ определяют относительную ширину судна и применяются в основном при оценке остойчивости, качки, ходкости и т.~д.

Особенности формы судовой поверхности достаточно хорошо характеризуются безразмерными коэффициентами полноты теоретического чертежа, причем эти коэффициенты могут вычисляться как для всего корпуса, так и для отдельных его частей, например, отдельно для кормовой или носовой.

\section{Коэффициенты полноты корпуса судна}

Из всех возможных безразмерных характеристик формы корпуса судна наибольшее применение для суждения о плавучести, остойчивости, ходкости, мореходности и управляемости получили следующие коэффициенты.

\textbf{Коэффициент общей полноты}, или коэффициент полноты водоизмещения $\delta$ --- отношение объема погруженной части корпуса судна $V$ к объему параллелепипеда со сторонами $L$, $B$, $T$, где $L$, $B$ --- длина и ширина действующей ватерлинии, $T$ --- осадка по эту ватерлинию.

\begin{equation}
    \delta = \frac{V}{L B T}.
\end{equation}

Действующей ватерлинией будет называться сечение судна поверхностью спокойной воды.
В общем случае действующая ватерлиния может не совпадать с ватерлиниями теоретического чертежа, и ее положение определяется посадкой судна.

\textbf{Коэффициент полноты ватерлинии} $\alpha$ --- отношение площади ватерлинии $S$ к площади прямоугольника со сторонами, равными длине и ширине этой ватерлинии

\begin{equation}
    \alpha = \frac{S}{L B}.
\end{equation}

\textbf{Коэффициент полноты площади шпангоута} $\beta$ --- отношение площади шпангоута по конкретную ватерлинию $\Omega$ к площади прямоугольника со сторонами, равными ширине и углублению (осадке) этого шпангоута по рассматриваемую ватерлинию

\begin{equation}
    \beta = \frac{\Omega}{B T}.
\end{equation}

Коэффициенты полноты корпуса судна $\delta$, $\beta$, $\alpha$ носят название основных.

\begin{figure}[h!]
    \centerline{\incfig{figure_01 - коэффициенты полноты}}
    \caption{К определению основных коэффициентов полноты}
\end{figure}

Нижеследующие коэффициенты называют производными от основных.

\textbf{Коэффициент продольной полноты} $\phi$ --- отношение объема погруженной части корпуса судна $V$ к объему цилиндра, площадь основания которого равна площади мидель-шпангоута, погруженного по рассматриваемую ватерлинию, а высота — длине судна по этой ватерлинии

\begin{equation}
    \phi = \frac{V}{\Omega \cdot L} = \frac{\delta L B T}{\beta B T L} = \frac{\delta}{\beta}.
\end{equation}

\textbf{Коэффициент вертикальной полноты} $\chi$ --- отношение объема погруженной части корпуса $V$ к объему цилиндра, площадь основания которого равна площади ватерлинии равновесия, а высота --- осадке судна по эту ватерлинию

\begin{equation}
    \chi = \frac{V}{S T} = \frac{\delta L B T}{\alpha L B T} = \frac{\delta}{\alpha}.
\end{equation}

% Соотношения главных размерений и коэффициентов полноты
\begin{table}[h]
    \caption{Соотношения главных размерений и коэффициентов полноты}
    \label{tab:ships}
    \centering
    \small
    \begin{tabular}{ | l | c | c | c | c | c | }
        \hline
        Типы судов & $L/B$ & $B/T$ & $\delta$ & $\alpha$ & $\beta$ \\ [1ex]
        \hline
        Пассажирские & 6,0--9,0 & 2,5--3,5 & 0,55--0,75 & 0,70--0,85 & 0,92--0,98 \\ [1ex]
        Грузовые     & 6,0--8,0 & 2,0--3,5 & 0,65--0,82 & 0,75--0,85 & 0,93--0,99 \\ [1ex]
        Танкеры      & 6,0--8,00 & 2,0--3,0 & 0,70--0,85 & 0,75--0,85 & 0,980--0,998 \\ [1ex]
        Ледоколы & 3,5--5,5 & 2,0--3,5 & 0,45--0,60 & 0,70--0,80 & 0,75--0,85 \\ [1ex]
        Тральщики & 6,4--7,5 & 3,5--4,3 & 0,50--0,60 & 0,65--0,80 & 0,75--0,95 \\ [1ex]
        \hline
    \end{tabular}
\end{table}

В таблице \ref{tab:ships} приведены характерные значения соотношений главных размерений и коэффициентов полноты по КВЛ для некоторых типов судов.

Помимо величин, приведенных в таблице, в расчетах теории корабля используются также коэффициент полноты диаметральной плоскости $\sigma$, представляющий собой отношение площади погруженной части диаметральной плоскости $A_L$ к площади прямоугольника со сторонами, равными длине судна по ватерлинии и осадке судна по эту ватерлинию

\begin{equation}
    \sigma = \frac{A_L}{L T}.
\end{equation}

\section{Расчеты и построение теоретического чертежа}

Расчеты по статике судов выполняются с использованием информации о поверхности корпуса судна.
Наиболее распространенным представлением такой информации является теоретический чертеж.
В настоящем пособии для вычерчивания теоретического чертежа используются исходные данные в безразмерном виде, приведенные в Приложении.

При составлении таблиц использованы следующие обозначения:

$\bar{y} = \frac{y}{B/2}$ --- Ординаты шпангоутов на соответствующих ватерлиниях;

$\text{П} = \frac{y_П}{B/2}$ --- Ординаты линии борта главной палубы;

$\bar{z} = \frac{z_П}{T}$ --- Аппликаты линии борта главной палубы;

$\bar{z}_1 = \frac{z_1}{T}$ --- Аппликаты шпангоутов на первом батоксе;

$\bar{z}_2 = \frac{z_2}{T}$ --- Аппликаты шпангоутов на втором батоксе;

$\bar{z}_{\text{Ф}} = \frac{z_{\text{Ф}}}{T}$ --- Аппликата точки пересечения форштевня с главной палубой;

$\bar{z}_{\text{А}} = \frac{z_{\text{А}}}{T}$ --- Аппликата точки пересечения ахтерштевня с главной палубой;

$\bar{x}_{\text{Ф}} = \frac{x_{\text{Ф}}}{\Delta L}$ --- Абсцисса точки пересечения форштевня с главной палубой относительно нулевого шпангоута;

$\bar{x}_{\text{А}} = \frac{x_{\text{А}}}{\Delta L}$ --- Абсцисса точки пересечения форштевня с главной палубой относительно десятого шпангоута.

В таблицах приложения даны значения безразмерных абсцисс, ординат и аппликат корпусов различных типов.
На соответствующих рисунках показаны корпуса этих судов, вычерченные в традиционной форме.

Для проведения расчетов по теории корабля необходимо вычертить корпус теоретического чертежа с некоторыми дополнительными построениями.

Исходными данными для выполнения расчетов является вариант таблицы безразмерных ординат и значения главных размерений (длина между перпендикулярами $L_{\text{ПП}}$, отношение длины к ширине $L/B$ и отношение ширины к осадке по КВЛ $B/T$).

Для вычерчивания теоретического чертежа необходимо привести таблицу безразмерных ординат к таблице размерных.
Выше показано, что обезразмеривание выполнено по величине полушироты $B/2$, осадки по КВЛ $T$ и теоретической шпации $\Delta L$, принятой в настоящих расчетах $\Delta L = L_{\text{ПП}}/10$.
Поэтому для получения размерных величин, приведенные в таблице величины следует умножить на соответствующие значения.

В таблице значения безразмерных ординат приведены на пересечении нескольких ватерлиний и шпангоутов.
Ватерлинии распределены равномерно, конструктивная ватерлиния обозначена на рисунке соответствующим сокращением (КВЛ).
Поэтому шаг по ватерлиниям считается $\Delta T = T / n$, где $n$ — порядковый номер ватерлинии, обозначенной КВЛ.

При расчетах статики шпангоуты должны быть построены до главной (верхней) палубы.
В таблице исходных данных кроме ординат приведены также ординаты и аппликаты бортовой линии палубы.

Расчеты статики выполняются и при больших наклонениях, поэтому на вспомогательном теоретическом чертеже должны быть построены линии палубы для каждого шпангоута.
Если не указано другого, палубная линия принимается выпуклой (имеется погибь бимсов).

Погибь бимсов принимается равной
\begin{equation}
    f = 1/25 \cdot y_{\text{П}},
\end{equation}
где $y_{\text{П}}$ — ордината верхней палубы на данном шпангоуте.

Для выполнения расчетов статики в отступлении от традиционного представления шпангоуты на теоретическом чертеже вычерчиваются на оба борта.
При этом шпангоуты в нос от мидель–шпангоута вычерчиваются сплошными линиями, а в корму — штриховыми.

Во времена, когда для проектирования судов не использовались компьютеры и САПР (системы автоматизированного проектирования), теоретические чертежи корпуса выполнялись на листах бумаги в соответствующем масштабе.
Проектирование с применением САПР позволяет выполнять моделирование (в том числе подготовку чертежей) в натурном размере (масштабе 1:1) и только отчетные документы выполнять в стандартных масштабах.
Таким образом, информация о судне отделена от формы ее представления.
В настоящем учебном пособии предполагается, что для вычерчивания теоретического чертежа корпуса используется та или иная система автоматизированного проектирования, поэтому все построения должны выполняться в натурном размере.
Такой подход исключает возникновение ошибки при масштабировании снятых с чертежа значений для использования в расчетах.
Если студенты вычерчивают теоретический чертеж на бумаге (или «миллиметровке»), то выполнять это следует, конечно же, в масштабе в соответствии с ЕСКД.
Следует также обратить внимание, что в практике проектирования судов для измерения линейных размеров корпуса принято использовать миллиметры, тогда как в расчетах статики используются метры.

При этом надо помнить, что в статике корабля применяется следующая система координат.
Начало системы располагается в точке $O$ --- точке пересечения трех основных плоскостей: диаметральной, мидель-шпангоута и основной.
Ось $X$ совпадает с линией пересечения диаметральной и основной плоскостей, положительное направление --- в нос.
Таким образом для точек, расположенных в нос от мидель-шпангоута, абсциссы положительны, для расположенных в корму --- отрицательны.
Ось $Y$ --- линия пересечения плоскостей мидель-шпангоута и основной, положительное направление на правый борт.
Таким образом, ординаты точек правой половины судна будут положительны, левой --- отрицательны.
Ось $Z$ --- линия пересечения плоскостей мидель-шпангоута и диаметральной, положительное направление вверх.
Таким образом, только для точек корпуса, расположенных ниже основной плоскости, аппликаты могут быть отрицательными.

% Построение линий шпангоутов
\begin{figure}[h!]
    \centerline{\incfig{figure_02 - 1}}
    \caption{Построение линий шпангоутов}
    \label{fig:lines-curves}
\end{figure}


На рисунке \ref{fig:lines-curves} изображено построение линий шпангоутов.
Сплошными кругами отмечены положения ординат на ватерлиниях.
Белыми кругами с серой кромкой --- положения аппликат линии штевней и батоксов.
Белый круг с черной кромкой отображает положение бортовой линии верхней палубы.
Следует проводить гладкую линию, проходящую через все точки (лучше всего использовать сплайн).
Если ординаты шпангоута имеют одинаковые значения на нескольких ватерлиниях (другими словами, у судна имеется вертикальный борт), то этот участок шпангоута лучше построить отрезком прямой линии.

Для того чтобы построить палубную линию, нужно нанести точку бортовой линии палубы  по ординате и аппликате из таблицы, затем посчитать и нанести в диаметральной плоскости аппликату палубной линии

\begin{equation}
    z = z_{\text{П}} + f = z_{\text{П}} + 1/25 \cdot y_{\text{П}}.
\end{equation}

Так как для расчетов статики требуется построить шпангоуты на оба борта, точка бортовой линии палубы откладывается зеркально на противоположный борт.
Через три точки проводится дуга окружности.

\begin{figure}[h!]
    \centerline{\incfig{figure_02 - 2}}
    \caption{Построение линий палубы}
    \label{fig:lines-curves-deck}
\end{figure}

\begin{figure}[h!]
    \centerline{\incfig{figure_02 - 3}}
    \caption{Вспомогательный теоретический чертеж}
    \caption*{\textit{Номерами отмечены шпангоуты, построенные на оба борта в отличие от классического представления теоретического чертежа.}}
    \label{fig:lines}
\end{figure}

\section{Методы численного интегрирования}

Большая часть вычислений статики корабля связана с вычислением интегралов --- объемных и по поверхности.
В связи с тем, что форма корпуса судна в общем виде не может быть задана аналитически, основной способ вычисления интегралов --- численный.
Ниже рассмотрено численное вычисление определенных интегралов и интегралов с переменным верхним пределом.

\subsection{Вычисление определенного интеграла}

Геометрически определенный интеграл можно представить как площадь фигуры, ограниченной кривой (заданной подынтегральным выражением), осью абсциссы, и вертикальными прямым, проходящими через нижний и верхний пределы интегрирования.

\begin{figure}[h!]
    \centerline{\incfig{figure_03 - 1}}
    \caption{Геометрический смысл определенного интеграла}
    \caption*{\textit{Определенный интеграл представляет собой площадь фигуры, ограниченной кривой, осью $X$ и вертикальными прямыми проходящими через точки начала и конца интегрирования --- $x_1$ и $x_2$.}}
\end{figure}

При этом кривая может представлять из себя зависимость линейной характеристики (например, ординаты ватерлинии), тогда интеграл представляет собой площадь, или зависимость площадей (например, зависимость площади ватерлинии от осадки), тогда интеграл представляет собой объем.
Иными словами, размерность определенного интеграла выражается произведением размерности подынтегральной функции и размерности переменной интегрирования.

Если $x$ в м, а $f(x)$ в м$^2$, тогда $\int f(x) dx $ --- в м$^3$.

Если $x$ в м, а $f(x)$ в м, тогда $\int f(x) dx $ --- в м$^2$.

Если $x$ в рад, а $f(x)$ в м, тогда $\int f(x) dx $ --- в рад $\cdot$ м.


Правило трапеции, используемое в настоящем учебном пособии, является одним из самых распространенных способов численного интегрирования.
Правило трапеций заключается в представлении площади под кривой в виде суммы площадей трапеций, построенных по нескольким точками, лежащим на этой кривой (\ref{fig:trapezoids}).

\begin{figure}[h!]
    \centerline{\incfig{figure_03 - 2}}
    \caption{Правило трапеций}
    \caption*{\textit{Линия представляет интегрируемую функцию на участке от $x_0$ до $x_3$.
    Значения функции в точках соединены ломаной линией, проведенной штриховой линией.
    Тогда площадь под кривой с определенной погрешностью представляется как площадь под ломаной линией.}}
    \label{fig:trapezoids}
\end{figure}

Обычно, кривую можно определить для нескольких значений переменной интегрирования, как с постоянным, так и с переменным шагом.
При использовании постоянного шага происходит значительное упрощение формулы определения.
Если значения подынтегрального выражения (функции $f(x)$) определены для значений $x_0$, $x_1$, $x_2$ и так далее, то в общем виде формула будет иметь вид

\begin{equation}
    \int f(x) dx = 0{,}5 (f_0 + f_1)(x_1 - x_0) + 0{,}5(f_1 + f_2)(x_2 - x_1)+ \dots
\end{equation}

Если в формуле выше принять, что шаг абсцисс постоянный и равен $\Delta x$, то формула приобретает вид

\begin{equation}
    \int f(x) dx = \Delta x (f_0 + f_1 + f_2 + \dots + f_n) - 0{,}5(f_0 + f_n).
\end{equation}

Из формулы выше видно, что по правилу трапеций определенный интеграл можно представить как сумму всех значений подынтегральной функции за вычетом полусуммы крайних значений, перемноженную на шаг по оси абсцисс.

В случае непостоянного шага, упрощения сделать не получится, и шаг требуется определять для каждой трапеции (при выполнении расчетов в табличном редакторе это приводит к появлению дополнительных столбцов в расчетных таблицах).

Повышения точности интегрирования можно добиться уменьшением шага (в пределе такая сумма переходит в интеграл) или переходом на способ численного интегрирования более высокого порядка (например, метод Симпсона или Чебышева).

\begin{figure}[h!]
    \centerline{\incfig{figure_06-1}}
    \caption{Учет истинного притыкания ватерлинии к ДП}
    \caption*{\textit{ В средней части приведен случай, когда судно имеет транцевую корму и ватерлиния оканчивается не пересечением с ДП.}}
\end{figure}

В практике судостроения принято разбивать поверхность корпуса и кривые, описывающие его, с постоянным шагом по теоретическим или практическим шпангоутам.
Рассмотрим для примера кривую ватерлинии (ВЛ), заданную значениями ординат по теоретическим шпангоутам.

Однако, нельзя гарантировать, что данная ватерлиния притыкается к диаметральной плоскости точно плоскости теоретического шпангоута.

А значит, что при выполнении численного интегрирования (например, для определения площади ватерлинии) значение площади будет отличаться от истинного на величину площадей треугольников, образованных точками притыкания ватерлинии к диаметральной плоскости и ординатами ватерлинии на соседних теоретических шпангоутах.

Площади треугольников можно представить выражениями

\begin{equation}
    \begin{split}
    A_{\text{Ф}} = 0{,}5 \Delta L_{\text{Ф}} f_{\text{Н}}, \\
    A_{\text{А}} = 0{,}5 \Delta L_{\text{А}} f_{\text{К}}.
    \end{split}
\end{equation}

Для того чтобы учесть при вычислении определенных интегралов несовпадение концевых точек подынтегральных кривых с теоретическими (или практическим шпангоутами) выполним следующие преобразования формулы правила трапеций

\begin{equation}
    \begin{split}
    \int f(x) dx = 0{,}5 \Delta L_{\text{Ф}} f_0 + 0{,}5 \Delta x (f_0 + f_1) + \dots \\
    + 0{,}5 \Delta x (f_{n-1} - f_n) + 0{,}5 \Delta L_{\text{А}} f_n.
    \end{split}
\end{equation}

Приведя общие члены, можно получить следующее выражение

\begin{equation}
    \int f(x) dx = \Delta L (K_{\text{Ф}} f_{\text{Н}} + f_{\text{Н}+1} + \dots + f_{\text{К}-1} + K_{\text{А}} f_{\text{К}}).
\end{equation}

В выражении выше коэффициенты вычисляются как

\begin{equation}
    \begin{split}
    K_{\text{Ф}} = 0{,}5 (1 + \Delta L_{\text{Ф}} / \Delta L), \\
    K_{\text{А}} = 0{,}5 (1 + \Delta L_{\text{А}} / \Delta L).
    \end{split}
\end{equation}

В случае наличия у судна транцевой кормы (транец --- плоский наклонный или вертикальный участок кормовой части судна), формулы для учета истинной точки окончания ватерлинии отличаются от вышеприведенных.
В этом случае дополнительная площадь имеет вид не треугольника, а трапеции.

\begin{equation}
    \begin{split}
    A_{\text{Ф}} = 0{,}5 \Delta L_{\text{Ф}} f_{\text{Н}}, \\
    A_{\text{А}} = 0{,}5 \Delta L_{\text{А}} \cdot (f_{\text{К}} + f_{\text{Т}}).
    \end{split}
\end{equation}

Тогда

\begin{equation}
    \begin{split}
    \int f(x) dx = 0{,}5 \Delta L_{\text{Ф}} f_0 + 0{,}5 \Delta x (f_0 + f_1) + \dots \\
    + 0{,}5 \Delta x (f_{n-1} - f_n) + 0{,}5 \Delta L_{\text{А}} (f_n + f_{\text{Т}}).
    \end{split}
\end{equation}

Приведя общие члены, можно получить следующее выражение

\begin{equation}
    \int f(x) dx = \Delta L (K_{\text{Ф}} f_{\text{Н}} + f_{\text{Н}+1} + \dots + f_{\text{К}-1} + K_{\text{А}} f_{\text{К}}).
\end{equation}

В выражении выше коэффициенты вычисляются как

\begin{equation}
    \begin{split}
    K_{\text{Ф}} = 0{,}5 (1 + \Delta L_{\text{Ф}}/ \Delta L), \\
    K_{\text{А}} = 0{,}5 (1 + \Delta L_{\text{А}} / \Delta L \cdot (1 + f_{\text{Т}} / f_{\text{А}})).
    \end{split}
\end{equation}

В случае наличия транцевой кормы в варианте задания будет особое об этом указание.

Рассмотрим на примере вычисление коэффициентов.

% Схема расчета коэффициентов $K_{\text{Ф}}$ и $K_{\text{А}}$ в табличном редакторе
\begin{figure}
    \footnotesize
    \centering
    \begin{tikzpicture}
        \matrix (mat) [table,
        row 1/.style={nodes={fill=gray!10,align=center,font=\segoe}},
        column 1/.style={nodes={fill=gray!10,align=center,font=\segoe}},
        column 1/.append style={nodes={text width=0.5cm}},
        column 2/.append style={nodes={text width=0.7cm}},
        column 3/.append style={nodes={text width=0.7cm}},
        column 4/.append style={nodes={text width=0.7cm}},
        column 5/.append style={nodes={text width=0.7cm}},
        column 6/.append style={nodes={text width=0.7cm}},
        column 7/.append style={nodes={text width=3.0cm}},
        column 8/.append style={nodes={text width=3.0cm}}
        ] {
        & A  & B & C & D & E & F & G \\
        1 & 0 & $\bar{x}_{\text{Ф}0}$ & $\bar{x}_{\text{А}0}$ & $i_{\text{Н}0}$ & $i_{\text{К}0}$ & =0,5*(1 + B1+D1) & =0,5*(1 + 10-C1-E1)\\
        2 & 1 & $\bar{x}_{\text{Ф}1}$ & $\bar{x}_{\text{А}1}$ & $i_{\text{Н}1}$ & $i_{\text{К}1}$ & =0,5*(1 + B2+D2) & =0,5*(1 + 10-C2-E2)\\
        3 & 2 & $\bar{x}_{\text{Ф}2}$ & $\bar{x}_{\text{А}2}$ & $i_{\text{Н}2}$ & $i_{\text{К}2}$ & =0,5*(1 + B3+D3) & =0,5*(1 + 10-C3-E3)\\
        4 & 3 & $\bar{x}_{\text{Ф}3}$ & $\bar{x}_{\text{А}3}$ & $i_{\text{Н}3}$ & $i_{\text{К}3}$ & =0,5*(1 + B4+D4) & =0,5*(1 + 10-C4-E4)\\
        5 & 4 & $\bar{x}_{\text{Ф}4}$ & $\bar{x}_{\text{А}4}$ & $i_{\text{Н}4}$ & $i_{\text{К}4}$ & =0,5*(1 + B5+D5) & =0,5*(1 + 10-C5-E5)\\
        6 & 5 & $\bar{x}_{\text{Ф}5}$ & $\bar{x}_{\text{А}5}$ & $i_{\text{Н}5}$ & $i_{\text{К}5}$ & =0,5*(1 + B6+D6) & =0,5*(1 + 10-C6-E6)\\
        7 & 6 & $\bar{x}_{\text{Ф}6}$ & $\bar{x}_{\text{А}6}$ & $i_{\text{Н}6}$ & $i_{\text{К}6}$ & =0,5*(1 + B7+D7) & =0,5*(1 + 10-C7-E7)\\
        8 & 7 & $\bar{x}_{\text{Ф}7}$ & $\bar{x}_{\text{А}7}$ & $i_{\text{Н}7}$ & $i_{\text{К}7}$ & =0,5*(1 + B8+D8) & =0,5*(1 + 10-C8-E8)\\
        };
    \end{tikzpicture}
    \caption{Схема расчета коэффициентов $K_{\text{Ф}}$ и $K_{\text{А}}$ в табличном редакторе}
    \caption*{\textit{В столбце $A$ --- номер шпангоута;
    в столбцах $B$ и $C$ --- относительные абсциссы штевней;
    в столбцах $D$ и $E$ --- номера шпангоутов, на которых оканчивается соответствующая ватерлиния;
    в столбцах $F$ и $G$ --- соответственно коэффициенты $K_{\text{Ф}}$ и $K_{\text{А}}$}}
    \label{fig:K-values}
\end{figure}

Величины $\bar{x}_{\text{Ф}}$ и $\bar{x}_{\text{А}}$ приведены от нулевого и десятого шпангоутов соответственно.
При этом положительные значение в нос, а отрицательные --- в корму.
Для расчета коэффициентов влияния необходимо вычислить расстояние от крайних точек на каждой ватерлинии до ближайшего шпангоута, то есть шпангоута, для которого в таблице ординат приведено значение.
Если в таблице для ватерлинии стоит знак «---», это означает, что ватерлиния не пересекает этот шпангоут.
На рисунке выше приведены номера крайних носовых и кормовых шпангоутов для каждой ватерлинии.

Формулы для вычисления расстояния от крайнего шпангоута до точки пересечения ватерлинии с ДП

\begin{equation}
    \begin{split}
    \Delta L_{\text{Ф}} = \Delta L (\bar{x}_{\text{Ф}} + i_{\text{Ф}}),\\
     \Delta L_{\text{А}} = \Delta L (10 - i_{\text{А}} - \bar{x}_{\text{А}}).
    \end{split}
\end{equation}

\subsection{Вычисление интеграла с переменным верхним пределом}

Кроме вычисления определенного интеграла в задачах статики выполняются вычисления интеграла с переменным верхним пределом вида

\begin{equation}
    \int_{z_0}^z \phi (z) dz,
\end{equation}
где $z_0$ --- значение (координата) начала интегрирования.

В отличие от определенного интеграла, результатом вычисления которого является одно число, интеграл с переменным верхним пределом приводит к получению функции, зависящей от той же переменной, что и подынтегральная функция.

Результат вычисления интеграла с переменным верхним пределом --- зависимость площади под кривой между началом интегрирования $z_0$ и текущим значением переменной $z$.

Рассмотрим способ вычисления интеграла с переменным верхним пределом на примере вычисления зависимости площади шпангоута от осадки $\Omega(z)$.
На рисунке \ref{fig:station-area-calculation} зависимость площади шпангоута от осадки приведена в столбце $D$.

\begin{equation}
    \Omega (z) = \int_{z_0}^z y(z) dz
\end{equation}

\begin{figure}[h!]
    \centerline{\incfig{figure_04 - 2}}
    \caption{К расчету интеграла с переменным верхним пределом}
    \caption*{\textit{На рисунке слева представлена схема расчета площади обычного шпангоута, когда нижняя точка кривой шпангоута не лежит на нулевой ватерлинии.
    В расчете используется коэффициент $K_0$, который учитывает, что начало интегрирования не совпадает с ватерлинией.
    На рисунке справа --- случай, когда необходимо использовать приведенное значение первой ординаты $y_0^*$}}
    \label{fig:station-area}
\end{figure}

% Схема расчета зависимости площади шпангоута от осадки в табличном редакторе
\begin{figure}
    \footnotesize
    \centering
        \begin{tikzpicture}
            \matrix (mat) [table,
            row 1/.style={nodes={fill=gray!10,align=center,font=\segoe}},
            column 1/.style={nodes={fill=gray!10,align=center,font=\segoe}},
            column 1/.append style={nodes={text width=0.5cm}},
            column 2/.append style={nodes={text width=1.5cm}},
            column 3/.append style={nodes={text width=1.5cm}},
            column 4/.append style={nodes={text width=3cm}},
            column 5/.append style={nodes={text width=3cm}}
            ] {
            & A & B & C & D \\
            1 & 0 & --- & 0 & =C1*$\Delta T$\\
            2 & $z_1$ & $y_1$ & =C1+0,5*B2*$K_0$ & =C2*$\Delta T$\\
            3 & $z_2$ & $y_2$ & =C2+0,5*(B2+B3) & =C3*$\Delta T$\\
            4 & $z_3$ & $y_3$ & =C3+0,5*(B3+B4) & =C4*$\Delta T$\\
            5 & $z_4$ & $y_4$ & =C4+0,5*(B4+B5) & =C5*$\Delta T$\\
            6 & $z_5$ & $y_5$ & =C5+0,5*(B5+B6) & =C6*$\Delta T$\\
            7 & $z_6$ & $y_6$ & =C6+0,5*(B6+B7) & =C7*$\Delta T$\\
            };
        \end{tikzpicture}
    \caption{Схема расчета зависимости площади шпангоута от осадки в табличном редакторе}
    \caption*{\textit{В столбце $B$ приведены ординаты шпангоута на соответствующих ватерлиниях;
    в столбце $C$ --- вспомогательные расчеты;
    в столбце $D$ --- зависимость площади шпангоута от осадки.}}
    \label{fig:station-area-calculation}
\end{figure}

В приведенном примере коэффициент, учитывающий, что начало интегрирования не совпадает с ватерлинией, вычисляется как

\begin{equation}
    K_0 = \Delta z / \Delta T.
\end{equation}

В случае, если интегрируемая кривая имеет большой перегиб между первой и второй координатами, для увеличения точности расчетов рекомендуется использовать приведенные значения первой ординаты.
Пример приведен справа на рисунке \ref{fig:station-area}.

Приведенная ордината $y_0^*$ определяется таким образом, чтобы площадь трапеции ($0{,}5(y_0^*+y_1)\Delta T$) равнялась истинной площади под кривой.
Приведенную ординату для каждого шпангоута (при необходимости) следует определить при построении вспомогательного теоретического чертежа.

Расчеты площади шпангоута при использовании приведенной ординаты отличаются в формулах для первой и второй ватерлиний.
Схема расчетов представлена на рисунке \ref{fig:station-area-2}.

%Схема расчета зависимости площади шпангоута от осадки в табличном редакторе. Приведенная ордината
\begin{figure}
    \footnotesize
    \centering
    \begin{tikzpicture}
        \matrix (mat) [table,
        row 1/.style={nodes={fill=gray!10,align=center,font=\segoe}},
        column 1/.style={nodes={fill=gray!10,align=center,font=\segoe}},
        column 1/.append style={nodes={text width=0.5cm}},
        column 2/.append style={nodes={text width=1.5cm}},
        column 3/.append style={nodes={text width=1.5cm}},
        column 4/.append style={nodes={text width=3cm}},
        column 5/.append style={nodes={text width=3cm}}
        ] {
        & A  & B & C & D \\
        1 & 0 & $y_0^*$ & 0 & =C1*$\Delta T$\\
        2 & $z_1$ & $y_1$ & =C1+0,5*(B1+B2) & =C2*$\Delta T$\\
        3 & $z_2$ & $y_2$ & =C2+0,5*(B2+B3) & =C3*$\Delta T$\\
        4 & $z_3$ & $y_3$ & =C3+0,5*(B3+B4) & =C4*$\Delta T$\\
        5 & $z_4$ & $y_4$ & =C4+0,5*(B4+B5) & =C5*$\Delta T$\\
        6 & $z_5$ & $y_5$ & =C5+0,5*(B5+B6) & =C6*$\Delta T$\\
        7 & $z_6$ & $y_6$ & =C6+0,5*(B6+B7) & =C7*$\Delta T$\\
        };
    \end{tikzpicture}
    \caption{Схема расчета зависимости площади шпангоута от осадки в табличном редакторе. Приведенная ордината}
    \caption*{\textit{В столбце $B$ приведены ординаты шпангоута на соответствующих ватерлиниях;
    в столбце $C$ --- вспомогательные расчеты;
    в столбце $D$ --- зависимость площади шпангоута от осадки.}}
    \label{fig:station-area-2}
\end{figure}

\section{Рекомендации по оформлению работы}

Табличный редактор (типа \textit{Microsoft Excel} или \textit{LibreOffice Calc}) является наиболее подходящей программой для выполнения расчетов курсового проекта и домашних заданий. 
Метод трапеций, описанный в настоящем учебном пособии, конечно, может использоваться и для ручного расчета с использованием калькулятора, однако, такой подход выходит за рамки целей пособия.

В соответствующих разделах методических указаний приведены требуемые для выполнения расчетов формулы и способы их реализации в табличном редакторе.

При выполнении курсового проекта или домашней работы рекомендуется следовать приведенным способам организации расчетов в файлах табличного редактора.
Рекомендуется организовать страницы файла соответственно видам расчетных работ: \textit{Исходные данные}, \textit{Расчет кривых элементов теоретического чертежа}, \textit{Расчет масштаба Бонжана} и т.д.

Расчеты в табличном редакторе всегда выполняются с машинной точностью, и результаты умножения или деления нецелых чисел могут иметь большое количество цифр после запятой.
В инженерных расчетах принято ограничивать количество значащих цифр в числе в зависимости от точности выполняемых вычислений.
При выполнении расчетов по статике морских судов рекомендуется отображать не более трех значащих цифр в любом приводимом в расчете числе.

Пример:

длина $L=100{,}5 \, \text{м}$ --- три значащие цифры после первой: 0, 0 и 5;

объем $V=40210 \, \text{м}^3$ --- три значащие цифры после первой: 0, 2 и 1.

При этом использовать функцию округления (например, ОКРУГЛ) в табличном редакторе не следует.
Используйте настройки формата отображения чисел в ячейках.
Для этого на выделенной ячейке выберите \textit{Формат ячеек} $\rightarrow$ \textit{Число}, затем выберите \textit{Числовой формат} и настройте необходимое число десятичных знаков.

После выполнения всех расчетов в табличном редакторе необходимо оформить отчет, включающий описание исходных данных, и для каждого вида расчетных работ использованные формулы, расчеты в табличном виде и полученные результаты.

Структура отчета должна включать следующие разделы:

\textbf{Введение}, в котором приводят цели выполнения работы;

\textbf{Исходные данные}, в котором приводят все исходные данные, полученные студентом вместе с заданием на курсовой проект или домашнюю работу;

разделы, соответствующие расчетным работам, например, \textbf{Расчет кривых элементов теоретического чертежа}.

Отчет оформляется в соответствии с требованиями правил ЕСКД для текстовых документов.

Графики и диаграммы, строящиеся при оценке мореходных качеств судов, служат не столько в качестве иллюстрации результатов вычислений, сколько для определения значений величин.
В связи с этой особенностью графики и диаграммы должны строиться (или вычерчиваться) с точно определенным масштабом, облегчающим снятие с них значений графическим способом.
Так, например, на графике площадей одному сантиметру на графике может соответствовать $1\,\text{м}^2$, $2\,\text{м}^2$, $2{,}5\,\text{м}^2$, $4\,\text{м}^2$, $5\,\text{м}^2$.
Масштаб построения кривых должен выбираться из стандартных масштабов, приведенных в нормативных документах ЕСКД.

Выполнение расчетов по теории корабля в табличном редакторе связано с набором однотипных формул в достаточно большое количество ячеек.
Для того чтобы не тратить много времени и не вписывать формулы в каждую ячейку, следует воспользоваться возможностью табличного редактора по «распространению» формулы на несколько ячеек.
При написании формулы со ссылкой на другие ячейки листа использование знака \$ перед наименованием столбца или строки позволяет при «распространении» формулы не изменять «зафиксированную» строку или столбец.
Например, если в формуле ссылка не содержит знака \$, то при «распространении» формулы вертикально или горизонтально, ссылка будет смещаться соответственно.
Если в ссылке зафиксирован столбец, например, \$B12, то при перемещении горизонтально, ссылка не сместится.
Если в ссылке зафиксирована строка, например, B\$12, то при перемещении вертикально, ссылка не сместится.
И, наконец, если зафиксирован и столбец, и строка, например, \$B\$12, то ссылка не изменится при любом «распространении» формулы.
